\section{Introduction}

Generative modeling within the realm of computer vision has experienced remarkable progress in recent years. Leading contemporary advancements are deep learning methodologies, specifically Generative Adversarial Networks (GANs) and diffusion models. These approaches have demonstrated exceptional capabilities in generating high-quality and diverse images. Nonetheless, challenges remain, such as balancing image fidelity and diversity while incorporating class-specific details, presenting ongoing research issues. Often, existing methods demand substantial computational resources and exhibit limited adaptability across various image generation contexts.

Motivated by the recent successes of reinforcement learning in optimizing complex decision-making tasks \cite{lillicrap2015continuous}, there has been a growing interest in Reinforced Adaptive Diffusion Networks (RAD-Nets) for image synthesis. These techniques endeavor to enhance image generation by integrating reinforcement signals to dynamically guide the diffusion process. The image synthesis process involves two primary stages: initial image generation from noise using diffusion models, followed by an optimization phase guided by reinforcement feedback. Unlike conventional methods adhering to static, data-driven procedures, these models leverage real-time feedback for refining generative parameters \cite{ho2020denoising}. Existing research indicates that this novel paradigm improves adaptability and fidelity compared to traditional GAN-based frameworks.

Current research primarily concentrates on improving image quality and diversity in isolation, without an integrated approach to concurrently ensure these objectives with class fidelity. Our findings suggest that static optimization techniques are inadequate, particularly in scenarios requiring rapid adjustments of model parameters for diverse datasets. Traditional diffusion models often struggle with achieving high fidelity in class-specific image synthesis, primarily due to their dependence on fixed parameter settings throughout the generation process \cite{song2020score}.

This paper introduces Reinforced Adaptive Diffusion Networks (RAD-Nets), a novel generative approach that integrates reinforcement learning techniques with diffusion models to iteratively refine outputs based on feedback-driven optimization. During inference, RAD-Nets adjust parameters dynamically in response to reinforcement signals, effectively balancing image quality, diversity, and fidelity. Specifically, RAD-Nets employ a multi-objective optimization framework, utilizing a quality-focused reward system to consistently guide parameter adjustments towards enhanced generative performance.

Our empirical evaluations on benchmark datasets such as CIFAR-10 and CelebA demonstrate that RAD-Nets achieve superior metrics for image quality and diversity compared to traditional models like GANs and unoptimized diffusion models. Notably, RAD-Nets exhibit an improved ability to preserve class-specific fidelity without compromising diversity. To the best of our knowledge, this study provides the first evidence of reinforcement-enhanced diffusion models surpassing static generative approaches across multiple evaluative metrics.

\begin{itemize}
    \item We introduce Reinforced Adaptive Diffusion Networks (RAD-Nets), a novel generative framework that combines diffusion models with reinforcement learning to dynamically optimize image generation processes.
    \item Our methodology integrates a multi-objective optimization framework that addresses image quality, diversity, and class fidelity, simultaneously achieving superior generative outcomes.
    \item We furnish empirical evidence showing that RAD-Nets outperform existing generative models on standard benchmarks, enhancing both fidelity and diversity.
    \item We illustrate the extensibility of RAD-Nets in class-conditional image synthesis, achieving significant improvements in class-specific feature representation and clarity.
\end{itemize}
